% Section 10.4, from lectures/L10/appendix/c_exercises.md. The worked solutions are not
% printed; the aside says where they are.

\section{Exercises}
\label{c10:sec:exercises}

Eight, ending with the capstone. That one is the whole book in one program: predict every stage, solve every stage, and reconcile the two, row by row.

\begin{aside}
\textbf{How to check your work.} A \kindtag{Code} exercise is judged by the suite that ships with it: run \code{make test} at the root of the companion repository, with \code{AEL_DIR} pointing at your toolkit. A suite you have not started is reported as \code{SKIP} rather than as a failure, so the command is worth running from the first day. The hand calculations and designs are checked against the toolkit functions each one names.

\textbf{The solutions are not in this book.} They are published in full in the companion repository, in \file{lectures/L10/appendix}, including the plausible wrong answer where there is one. Read them once you have your own answers, including the ones you are unsure about, because being unsure and right is a different thing from being unsure and wrong and the solutions say which is which.
\end{aside}

% ------------------------------------------------------------------
\recall{What each stage is for}
\label{c10:ex:1}
\begin{enumerate}
  \item Name the four stages and say what each is responsible for.
  \item Two of them have no voltage gain. Which, and what are they for?
  \item Which stage decides the amplifier's input offset, and can a later stage improve it?
  \item Which node carries the compensation, and why that one?
\end{enumerate}

% ------------------------------------------------------------------
\recall{The loop}
\label{c10:ex:2}
\begin{enumerate}
  \item Write the closed-loop gain in terms of the open-loop gain and the feedback fraction.
  \item What is the loop gain of this amplifier at a closed-loop gain of 20?
  \item Feedback divides the output resistance by one plus the loop gain. Compute it, then say why the answer is not a real number.
  \item The open-loop gain varies by a factor of 14 across a beta spread. By how much does the closed-loop gain vary, and why is that the answer feedback exists to give?
\end{enumerate}

% ------------------------------------------------------------------
\handcalc{The budget}
\label{c10:ex:3}
The amplifier of \secref{c10:app:a}: 2 mA tail, 1 mA buffer, 1 mA voltage stage, 120 mA idle, 8 ohm load, $\beta = 50$.

\begin{enumerate}
  \item The unloaded gain of stages 1 and 3.
  \item The input resistance of stages 2, 3 and 4.
  \item Each stage loaded by the next, and the open-loop gain.
  \item The difference in decibels between the product of the unloaded gains and your answer to part 3. Where is most of it?
\end{enumerate}
\checkyourself{\code{budget}, \code{openLoopGain}}

% ------------------------------------------------------------------
\handcalc{The stages that do nothing}
\label{c10:ex:4}
\begin{enumerate}
  \item What would stage 1's gain be if it drove stage 3 directly? Give the open-loop gain that results.
  \item What would stage 3's gain be if the output stage were a single follower rather than a Darlington? Give the open-loop gain that results.
  \item State what each of the two no-gain stages is worth, in decibels.
  \item Both stages cost signal. State what each costs, and comment on the exchange rate.
\end{enumerate}

% ------------------------------------------------------------------
\design{The input offset that matters}
\label{c10:ex:5}
\begin{enumerate}
  \item What input offset voltage puts the open-loop output at a rail?
  \item Two ordinary transistors match their $V_{BE}$ to about 1 mV. How many times larger is that than your answer to part 1?
  \item State what this means for solving the amplifier open-loop.
  \item The closed-loop amplifier has an output offset. Compute it for 1 mV of input offset at a closed-loop gain of 20, and say what would reduce it.
\end{enumerate}

% ------------------------------------------------------------------
\design{The gain you cannot buy}
\label{c10:ex:6}
\begin{enumerate}
  \item Compute stage 1's gain at a tail current of 0.2 mA, 2 mA and 20 mA, with a mirror load in each case. Then do the same for stage 3 at 0.1 mA, 1 mA and 10 mA with a current-source load.
  \item Explain the result. Write the gain symbolically and show what cancels.
  \item So how \emph{do} you get more gain out of one stage? Compute what your answer gives at 1 mA.
  \item The compensation capacitor is 30 pF. Compute the amplifier's slew rate at the 2 mA tail, then the tail current that would give 10 V per microsecond, and say what that change costs in gain.
\end{enumerate}
\checkyourself{\code{mirrorGain}, \code{intrinsicEmitterResistance}, \code{cascodeOutputResistance}}

% ------------------------------------------------------------------
\codeexercise{PNP support, and the report}
\label{c10:ex:7}
Add \code{Polarity} to \code{addBjt} and stamp a PNP as an NPN with every voltage and current negated. Then implement \code{ael::report} to the specification in \secref{c10:sec:b6}.

\textbf{\code{budget} must compose the functions from Chapter~\ref{c7:ch:smallsignal} to Chapter~\ref{c9:ch:diffpair}, not restate their formulas.} One shipped test checks that its stage-1 figure agrees with \code{ael::diffpair::mirrorGain} to machine precision, which it can only do by calling it.

% ------------------------------------------------------------------
\crosscheck{The capstone}
\label{c10:ex:8}
\textbf{Predict the whole amplifier, solve the whole amplifier, and reconcile the two, row by row.}

\begin{enumerate}
  \item \textbf{Predict.} Run \code{budget} and record the four stage gains and the open-loop gain.
  \item \textbf{Solve, stage by stage.} Build each stage as its own netlist, bias it, perturb its input by a small voltage, and measure its output. Load each stage with the next stage's input resistance as a resistor, so that the solver is measuring the same quantity the closed form predicted.
  \item \textbf{Reconcile.} Print predicted against measured with the difference in per cent, one row per stage.
  \item \textbf{Then solve the whole amplifier open-loop}, with the four stages connected. Report what happens.
  \item \textbf{Then close the loop} with a feedback divider for a gain of 20, solve again, and report the closed-loop gain.
\end{enumerate}

\task{What to expect}
\textbf{Every row differs by a few per cent, and every difference has a name.} The closed forms take $V_T$ as 26 mV where the device computes $kT/q = 25.87$; they take $r_o$ as $V_A/I_C$ where differentiating the model gives $(V_A + V_{CE})/I_C$; they take $\beta$ as 50 where the transport model's incremental value differs slightly. \textbf{Account for every row.} A reader whose rows agree to six figures has run the same code twice and learned nothing.

\textbf{Part 4 will not converge on a sensible answer, and that is the correct result.} With an open-loop gain of a million and rails of 15 V, an input offset of \textbf{15 microvolts} saturates the output. Your solver will report the output at a rail, or fail to converge, and both are true statements about operational amplifiers. Do not tune your solver until it produces a mid-rail answer; there is not one.

\textbf{Part 5 will converge}, and the closed-loop gain should be within a per cent of 20. Notice what you have just demonstrated: the same circuit is unsolvable open-loop and well-behaved closed-loop, and nothing about the transistors changed.

\textbf{Finally, the exercise that is the book.} Repeat part 5 with $\beta$ at 20 and at 200. The open-loop gain moves by a factor of 14. Report what the closed-loop gain does, and say in one sentence what that is worth.
